%
% Informe Práctica 3 — Aprendizaje por Refuerzo en Pendulum-v1
% MAAAII · Curso 2025/2026
%
\documentclass[12pt,a4paper]{article}

% ── Paquetes ──────────────────────────────────────────────────────────────────
\usepackage[utf8]{inputenc}
\usepackage[T1]{fontenc}
\usepackage[spanish]{babel}
\usepackage{lmodern}
\usepackage{geometry}
\geometry{margin=2.5cm}
\usepackage{amsmath,amssymb}
\usepackage{booktabs}
\usepackage{graphicx}
\usepackage{float}
\usepackage{hyperref}
\hypersetup{colorlinks=true, linkcolor=blue, urlcolor=blue, citecolor=blue}
\usepackage{caption}
\usepackage{subcaption}
\usepackage{xcolor}
\usepackage{enumitem}
\usepackage{parskip}

% ── Metadatos ─────────────────────────────────────────────────────────────────
\title{%
  \textbf{Práctica 3: Aprendizaje por Refuerzo}\\[4pt]
  \large Modelos Avanzados de Aprendizaje Automático II\\
  Grado en Inteligencia Artificial · Curso 2025/2026
}
\author{Grupo: jgdans -- pdcobelo}
\date{Mayo 2026}

% ══════════════════════════════════════════════════════════════════════════════
\begin{document}
\maketitle
\tableofcontents
\newpage

% ══════════════════════════════════════════════════════════════════════════════
\section{Descripción del entorno y enfoque general}

El entorno \texttt{Pendulum-v1} de Gymnasium simula un péndulo simple al que se le
puede aplicar un par (torque) rotacional. El objetivo es mantener el péndulo en
posición vertical hacia arriba. El espacio de observación es continuo y tridimensional:
$(\cos\theta,\,\sin\theta,\,\dot\theta)$, donde $\theta$ es el ángulo respecto a la
posición vertical y $\dot\theta$ es la velocidad angular. La acción también es continua:
el torque aplicado $\tau \in [-2{,}2]$~N·m.

La función de recompensa es:
\[
  r(\theta,\dot\theta,\tau)
  = -\bigl(\theta^2 + 0{,}1\,\dot\theta^2 + 0{,}001\,\tau^2\bigr),
\]
siempre negativa, con máximo 0 cuando el péndulo está perfectamente vertical y estático.
Cada episodio tiene $T=200$ pasos con $\Delta t=0{,}05$~s, es decir exactamente
\textbf{10 segundos} de simulación, tal y como exige el enunciado.

Para aplicar algoritmos de RL tabular es necesario discretizar tanto el espacio de
estados como el de acciones, lo que se describe en la sección siguiente.

% ══════════════════════════════════════════════════════════════════════════════
\section{Pregunta 1: Discretización aplicada}

\subsection{Estado}

La observación $(\cos\theta,\sin\theta,\dot\theta)$ se comprime a dos variables
físicamente interpretables:

\begin{enumerate}
  \item \textbf{Ángulo} $\theta = \operatorname{atan2}(\sin\theta,\cos\theta)
        \in [-\pi,\pi]$, discretizado en $N_\theta=20$ intervalos iguales.
  \item \textbf{Velocidad angular} $\dot\theta\in[-8,8]$~rad/s
        (límites del entorno), discretizada en $N_{\dot\theta}=20$ intervalos.
\end{enumerate}

Recuperar $\theta$ mediante \texttt{atan2} elimina la discontinuidad en $\pm\pi$
sin necesidad de trabajar con tres dimensiones de estado.

\subsection{Acción}

El torque continuo $\tau\in[-2{,}2]$ se discretiza en $N_a=9$ valores equiespaciados:
$\{-2{,}0,\,-1{,}5,\,-1{,}0,\,-0{,}5,\,0{,}0,\,0{,}5,\,1{,}0,\,1{,}5,\,2{,}0\}$.

\subsection{Espacio total y justificación}

El espacio resultado tiene $20\times 20\times 9=3\,600$ pares estado-acción
(tabla~\ref{tab:discretizacion}).

\begin{table}[H]
  \centering
  \caption{Parámetros de discretización.}
  \label{tab:discretizacion}
  \begin{tabular}{llcc}
    \toprule
    Variable & Rango & Bins & Paso \\
    \midrule
    $\theta$ & $[-\pi,\pi]$ & 20 & $\approx 0{,}314$~rad \\
    $\dot\theta$ & $[-8,8]$~rad/s & 20 & $0{,}8$~rad/s \\
    $\tau$ & $[-2,2]$~N·m & 9 & $0{,}5$~N·m \\
    \bottomrule
  \end{tabular}
\end{table}

Con 20 bins por dimensión de estado se distinguen los estados cualitativamente
relevantes (péndulo arriba/abajo, girando rápido/despacio), sin aumentar el espacio
hasta el punto en que los algoritmos tabulares no converjan en 5\,000 episodios.
Discretizaciones más gruesas (p.ej.\ 10 bins) fusionarían estados con dinámica
distinta; más finas (p.ej.\ 50 bins) elevarían el espacio a $22\,500$ pares,
requiriendo muchos más episodios para visitar cada celda suficientemente.

% ══════════════════════════════════════════════════════════════════════════════
\section{Pregunta 2: Hiperparámetros}

\subsection{Valores utilizados}

Los mismos hiperparámetros se aplican a los tres algoritmos para que la comparativa
sea equitativa (tabla~\ref{tab:hparams}).

\begin{table}[H]
  \centering
  \caption{Hiperparámetros base de los tres algoritmos.}
  \label{tab:hparams}
  \begin{tabular}{clp{8.5cm}}
    \toprule
    Parámetro & Valor & Justificación \\
    \midrule
    $\gamma$ & 0{,}99 & Pendulum requiere mantener el equilibrio durante 200 pasos;
                        un $\gamma$ menor penalizaría en exceso los retornos tardíos. \\
    $\alpha$ & 0{,}10 & Tasa de aprendizaje conservadora y estable (ver ablación). \\
    $\varepsilon$ & 0{,}10 & Equilibrio entre exploración y explotación (ver ablación). \\
    Episodios & 5\,000 & Suficientes para observar convergencia con 3\,600 pares. \\
    \bottomrule
  \end{tabular}
\end{table}

\subsection{Selección de $\alpha$ mediante ablación}

Se entrenó Q-Learning durante 3\,000 episodios variando $\alpha$ con $\varepsilon=0{,}1$
fijo. La métrica es el retorno medio de los últimos 500 episodios de entrenamiento
(figura~\ref{fig:ablacion_alpha} y tabla~\ref{tab:ablacion_alpha}).

\begin{figure}[H]
  \centering
  \includegraphics[width=0.7\textwidth]{fig_ablacion_alpha.pdf}
  \caption{Ablación de $\alpha$: retorno medio en los últimos 500 episodios de
           entrenamiento (Q-Learning, $\varepsilon=0{,}1$, 3000 ep.).}
  \label{fig:ablacion_alpha}
\end{figure}

\begin{table}[H]
  \centering
  \caption{Valores de ablación de $\alpha$ ($\varepsilon=0{,}1$).}
  \label{tab:ablacion_alpha}
  \begin{tabular}{cc}
    \toprule
    $\alpha$ & Retorno (últ.\ 500 ep.) \\
    \midrule
    0{,}01 & $-1165{,}58$ \\
    0{,}05 & $-899{,}74$  \\
    \textbf{0{,}10} & $\mathbf{-730{,}03}$ \\
    0{,}20 & $-597{,}33$  \\
    0{,}50 & $-653{,}68$  \\
    \bottomrule
  \end{tabular}
\end{table}

$\alpha=0{,}20$ da el mejor retorno de entrenamiento en 3\,000 episodios, lo que indica
convergencia más rápida. Sin embargo, $\alpha=0{,}50$ ya empeora ($-653$), mostrando
inestabilidad por sobreactualización. $\alpha=0{,}10$ es la opción más conservadora:
converge algo más lento pero de forma estable, y sus actualizaciones tienen menor
varianza, lo que es beneficioso para políticas finales más consistentes a lo largo de
5\,000 episodios.

\subsection{Selección de $\varepsilon$ mediante ablación}

Se repitió el estudio variando $\varepsilon$ con $\alpha=0{,}1$ fijo
(figura~\ref{fig:ablacion_eps} y tabla~\ref{tab:ablacion_eps}).

\begin{figure}[H]
  \centering
  \includegraphics[width=0.7\textwidth]{fig_ablacion_eps.pdf}
  \caption{Ablación de $\varepsilon$: retorno medio en los últimos 500 episodios de
           entrenamiento (Q-Learning, $\alpha=0{,}1$, 3000 ep.).}
  \label{fig:ablacion_eps}
\end{figure}

\begin{table}[H]
  \centering
  \caption{Valores de ablación de $\varepsilon$ ($\alpha=0{,}1$).}
  \label{tab:ablacion_eps}
  \begin{tabular}{cc}
    \toprule
    $\varepsilon$ & Retorno (últ.\ 500 ep.) \\
    \midrule
    0{,}01 & $-514{,}47$ \\
    0{,}05 & $-589{,}35$ \\
    \textbf{0{,}10} & $\mathbf{-731{,}81}$ \\
    0{,}20 & $-826{,}19$ \\
    0{,}50 & $-961{,}22$ \\
    \bottomrule
  \end{tabular}
\end{table}

El retorno de entrenamiento crece monótonamente al reducir $\varepsilon$: con
$\varepsilon=0{,}01$ el agente apenas explora y se comporta casi como greedy durante
el entrenamiento, lo que infla artificialmente el retorno medio medido. Sin embargo,
explorar poco al inicio de Q-Learning puede dejar regiones del espacio de estados sin
visitar y llevar a mínimos locales. $\varepsilon=0{,}10$ ofrece suficiente exploración
para cubrir el espacio de estados en 5\,000 episodios. La variante con
\emph{$\varepsilon$ decreciente} ($0{,}10 \to 0{,}01$) combina lo mejor de ambos
extremos y se analiza en la sección~\ref{sec:decay}.

\subsection{Efecto del decaimiento de $\varepsilon$}
\label{sec:decay}

Se entrenó Q-Learning con $\varepsilon$ decreciente de forma lineal desde $0{,}10$
hasta $0{,}01$ a lo largo de los 5\,000 episodios. Los resultados muestran una mejora
notable en la política final:

\begin{center}
\begin{tabular}{lcc}
  \toprule
  Configuración & Retorno greedy (media) & Desv. est. \\
  \midrule
  Q-Learning $\varepsilon=0{,}1$ fijo  & $-304{,}69$ & $274{,}06$ \\
  Q-Learning $\varepsilon$ decreciente & $-223{,}20$ & $138{,}87$ \\
  \bottomrule
\end{tabular}
\end{center}

El decaimiento permite explorar ampliamente al inicio (mayor cobertura del espacio)
y explotar la política aprendida al final (mayor retorno greedy). Además reduce la
varianza, lo que indica una política más consistente.

% ══════════════════════════════════════════════════════════════════════════════
\section{Pregunta 3: Política final — mejor algoritmo determinista}

Se evaluó la política greedy de cada algoritmo durante 100 episodios con semilla fija.
Los resultados se recogen en la tabla~\ref{tab:greedy}.

\begin{table}[H]
  \centering
  \caption{Retorno greedy final (100 episodios, sin ruido, semilla única).}
  \label{tab:greedy}
  \begin{tabular}{lcc}
    \toprule
    Algoritmo & Retorno medio & Desv.\ est. \\
    \midrule
    Monte Carlo ($\varepsilon$-soft on-policy) & $-499{,}42$ & $146{,}94$ \\
    SARSA                                      & $-504{,}26$ & $441{,}58$ \\
    Q-Learning ($\varepsilon$ fijo)            & $-875{,}72$ & $303{,}98$ \\
    Q-Learning ($\varepsilon$ decreciente)     & $-223{,}20$ & $138{,}87$ \\
    \bottomrule
  \end{tabular}
\end{table}

Para confirmar la robustez estadística se repitió el entrenamiento con tres semillas
y se evaluó con 50 episodios greedy por semilla (figura~\ref{fig:multiseed} y
tabla~\ref{tab:multiseed}).

\begin{figure}[H]
  \centering
  \includegraphics[width=0.7\textwidth]{fig_multiseed.pdf}
  \caption{Evaluación multi-semilla (seeds 0, 1, 2): retorno greedy medio $\pm$ std.}
  \label{fig:multiseed}
\end{figure}

\begin{table}[H]
  \centering
  \caption{Evaluación multi-semilla (política greedy, 50 episodios/semilla).}
  \label{tab:multiseed}
  \begin{tabular}{lccc}
    \toprule
    Algoritmo & Media & Std & Por semilla (0, 1, 2) \\
    \midrule
    MC                          & $-314{,}93$ & $96{,}54$  & $-221{,}3$ / $-447{,}8$ / $-275{,}7$ \\
    SARSA                       & $-331{,}92$ & $11{,}09$  & $-342{,}3$ / $-316{,}5$ / $-336{,}9$ \\
    Q-Learning                  & $-485{,}78$ & $164{,}32$ & $-601{,}0$ / $-602{,}9$ / $-253{,}4$ \\
    QL + $\varepsilon$-decay    & $-330{,}43$ & $103{,}24$ & $-232{,}8$ / $-473{,}3$ / $-285{,}2$ \\
    \bottomrule
  \end{tabular}
\end{table}

\textbf{Análisis:} con una sola semilla \textbf{Q-Learning con $\varepsilon$ decreciente}
obtiene el mejor retorno greedy ($-223{,}20$), seguido de Monte Carlo ($-499{,}42$).
Con múltiples semillas, MC ($-314{,}93$) lidera, con SARSA ($-331{,}92$) y
QL+$\varepsilon$-decay ($-330{,}43$) prácticamente empatados.

El \textbf{bajo rendimiento de Q-Learning con $\varepsilon$ fijo} tiene dos causas:
\begin{enumerate}
  \item \emph{Exploración tardía:} con $\varepsilon=0{,}1$ constante el agente sigue
        explorando al final del entrenamiento, impidiendo que $Q$ converja hacia la
        política óptima.
  \item \emph{Optimismo del} $\max$: la actualización off-policy usa
        $\max_{a'}Q(s',a')$, sobreestimando el valor de estados poco visitados
        (inicializados a 0). Esto sesga las actualizaciones en un entorno donde
        todas las recompensas son negativas.
\end{enumerate}

Monte Carlo, on-policy y con retornos de episodio completo, estima $Q$ de forma
insesgada sin sobreestimación del $\max$, resultando en una política robusta a lo
largo de distintas semillas.

La figura~\ref{fig:heatmaps} muestra los mapas de calor de la política greedy de los
tres algoritmos entrenados con semilla~0. En cada celda se muestra el torque elegido
(rojo = torque positivo, azul = torque negativo) en función del ángulo $\theta$ y la
velocidad angular $\dot\theta$.

\begin{figure}[H]
  \centering
  \includegraphics[width=\textwidth]{fig_heatmaps.pdf}
  \caption{Política greedy aprendida por cada algoritmo: mapa de calor del torque
           elegido en función del estado $(\theta, \dot\theta)$.}
  \label{fig:heatmaps}
\end{figure}

MC y SARSA muestran patrones estructurados con zonas claras de torque positivo y
negativo simétricas respecto al eje central, lo que corresponde a una estrategia
coherente de control. La política de Q-Learning (con $\varepsilon$ fijo) muestra un
patrón más irregular, consistente con su peor rendimiento cuantitativo.

\textbf{Conclusión:} \textbf{Monte Carlo} obtiene la política determinista más
consistente en la evaluación multi-semilla. Q-Learning con $\varepsilon$ decreciente
es competitivo en la semilla~0 pero su std es elevada.

% ══════════════════════════════════════════════════════════════════════════════
\section{Pregunta 4: Rendimiento durante el aprendizaje}

La figura~\ref{fig:curvas} muestra las curvas de aprendizaje (retorno medio móvil de
100 episodios) durante el entrenamiento.

\begin{figure}[H]
  \centering
  \includegraphics[width=0.85\textwidth]{fig_curvas.pdf}
  \caption{Curvas de aprendizaje para MC, SARSA y Q-Learning ($\varepsilon$ fijo).}
  \label{fig:curvas}
\end{figure}

Los retornos medios móviles al final del entrenamiento (episodio~4900) son:

\begin{center}
\begin{tabular}{lc}
  \toprule
  Algoritmo & Retorno (media móvil, ep.\ 4900) \\
  \midrule
  Monte Carlo           & $-266{,}32$ \\
  SARSA                 & $-934{,}10$ \\
  Q-Learning fijo       & $-478{,}81$ \\
  Q-Learning decay      & $-226{,}90$ \\
  \bottomrule
\end{tabular}
\end{center}

\textbf{Análisis:}

\begin{itemize}
  \item \textbf{Monte Carlo} obtiene el mejor retorno de entrenamiento al final
        ($-266{,}32$). Aunque su señal tiene alta varianza (retornos completos de
        200 pasos), el promedio incremental permite convergencia estable.
  \item \textbf{Q-Learning con $\varepsilon$ decreciente} es muy competitivo
        ($-226{,}90$), confirmando que el decaimiento de $\varepsilon$ mejora
        tanto la política final como el rendimiento durante el aprendizaje.
  \item \textbf{Q-Learning fijo} tiene retorno de entrenamiento medio
        ($-478{,}81$) a pesar de su pobre política greedy, porque los episodios
        de entrenamiento incluyen exploración aleatoria.
  \item \textbf{SARSA} muestra el peor retorno de entrenamiento ($-934{,}10$),
        lo que puede reflejar inestabilidad en la estimación on-policy con
        $\varepsilon$ constante en este entorno de recompensas muy negativas.
\end{itemize}

\textbf{Conclusión:} \textbf{Monte Carlo} obtiene el mejor rendimiento como
algoritmo de control (mayor retorno medio durante el aprendizaje). Su carácter
on-policy y la estimación insesgada mediante retornos reales le permiten aprender
una política equilibrada desde las primeras iteraciones.

% ══════════════════════════════════════════════════════════════════════════════
\section{Pregunta 5: Robustez ante perturbaciones}

\subsection{Experimento}

Con probabilidad $p=0{,}05$, la acción elegida por el agente se reemplaza por
$[-2{,}0]$ o $[2{,}0]$ (torque máximo aleatorio). Los resultados se muestran en la
tabla~\ref{tab:ruido} y la figura~\ref{fig:robustez}.

\begin{table}[H]
  \centering
  \caption{Retorno greedy con y sin ruido ($p=0{,}05$).
           $\Delta=$ retorno con ruido $-$ retorno sin ruido.}
  \label{tab:ruido}
  \begin{tabular}{lrrr}
    \toprule
    Algoritmo & Sin ruido & Con ruido & $\Delta$ \\
    \midrule
    Monte Carlo   & $-499{,}42$ & $-518{,}48$ & $-19{,}05$ \\
    SARSA         & $-504{,}26$ & $-536{,}03$ & $-31{,}77$ \\
    Q-Learning    & $-875{,}72$ & $-831{,}69$ & $+44{,}03$ \\
    QL robusto    & $-409{,}19$ & $-462{,}37$ & $-53{,}18$ \\
    \bottomrule
  \end{tabular}
\end{table}

\begin{figure}[H]
  \centering
  \includegraphics[width=0.75\textwidth]{fig_robustez.pdf}
  \caption{Comparación de retorno greedy con y sin ruido para cada algoritmo.}
  \label{fig:robustez}
\end{figure}

\subsection{¿Qué algoritmo es más robusto?}

Atendiendo a la caída absoluta de rendimiento al introducir ruido, \textbf{Monte Carlo}
es el más robusto ($|\Delta|=19{,}05$), seguido de SARSA ($|\Delta|=31{,}77$).

El $\Delta>0$ de Q-Learning es aparentemente paradójico: el rendimiento \emph{mejora}
con ruido. La explicación es que la política greedy de Q-Learning con $\varepsilon$
fijo es de baja calidad ($-875{,}72$), y los torques extremos aleatorios ($\pm2$~N·m)
aportan de forma accidental la energía necesaria para elevar el péndulo, compensando
el defecto de la política. En otras palabras, el ruido ``ayuda'' a una política
deficiente; la política no es intrínsecamente robusta.

Comparando solo los algoritmos con buen rendimiento base,
\textbf{Monte Carlo es el más robusto}.

\subsection{Adaptación del aprendizaje: domain randomization}

Se reentrenó Q-Learning inyectando el mismo tipo de ruido ($p=0{,}05$) durante el
propio entrenamiento (\emph{domain randomization}):

\begin{center}
\begin{tabular}{lrrr}
  \toprule
  Variante & Sin ruido & Con ruido ($p=0{,}05$) & $\Delta$ \\
  \midrule
  Q-Learning original & $-875{,}72$ & $-831{,}69$ & $+44{,}03$ \\
  Q-Learning robusto  & $-409{,}19$ & $-462{,}37$ & $-53{,}18$ \\
  \bottomrule
\end{tabular}
\end{center}

El Q-Learning robusto mejora el retorno base en $\sim\!470$ puntos y pasa a tener
el mejor retorno con ruido entre todos los agentes evaluados ($-462{,}37$).

\textbf{¿Por qué funciona?} El ruido en entrenamiento actúa como
\textbf{exploración adicional dirigida}: fuerza visitas a estados de alta energía
(péndulo caído, velocidad angular extrema) y obliga al agente a aprender a
recuperarse. Esto amplía la cobertura del espacio de estados y genera una tabla $Q$
más densa en regiones críticas, explicando también la mejora en la política greedy
sin ruido.

\textbf{Otras estrategias posibles para mejorar la robustez:}
\begin{itemize}
  \item \textbf{Modelar la incertidumbre explícitamente:} tratar el ruido como
        parte del modelo de transición (RL robusto, Interval MDP).
  \item \textbf{Ensemble de políticas:} entrenar múltiples agentes con distintas
        realizaciones de ruido y combinar sus políticas mediante voto.
  \item \textbf{$\varepsilon$ decreciente combinado con domain randomization:}
        exploración amplia al inicio, ruido de perturbación realista durante todo
        el entrenamiento y explotación al final.
\end{itemize}

% ══════════════════════════════════════════════════════════════════════════════
\section{Conclusiones}

\begin{enumerate}
  \item La discretización $20\times20$ de estado y 9 acciones ($3\,600$ pares
        en total) ofrece resolución suficiente para distinguir la dinámica del péndulo
        y converge en $5\,000$ episodios.
  \item Los hiperparámetros $\alpha=0{,}1$ y $\varepsilon=0{,}1$ se justifican
        mediante ablación: $\alpha$ mayor converge más rápido pero con mayor
        varianza; $\varepsilon$ menor infla el retorno de entrenamiento sin garantizar
        mayor calidad de política. La variante con $\varepsilon$ decreciente combina
        exploración amplia e initial con explotación final.
  \item En evaluación de política final, \textbf{Monte Carlo} es el más consistente
        a lo largo de múltiples semillas. Q-Learning con $\varepsilon$ decreciente
        es el mejor en semilla única.
  \item \textbf{Monte Carlo} obtiene el mejor rendimiento durante el aprendizaje
        (mayor retorno medio de entrenamiento). Q-Learning con $\varepsilon$
        decreciente es competitivo.
  \item \textbf{Monte Carlo} es el más robusto ante perturbaciones (menor caída
        relativa). El domain randomization en Q-Learning lo convierte en el agente
        con mayor retorno absoluto bajo ruido.
\end{enumerate}

\end{document}
